\documentclass{article}
\usepackage{amsmath,amssymb}
\usepackage{booktabs}
\usepackage{geometry}
\usepackage{enumitem}
\geometry{margin=1in}

\title{Notes on the proposal: connecting papers, data, and the search for the
fundamental tension}
\author{Companion notes to \texttt{proposal\_eee\_draft.tex}}
\date{}

\begin{document}
\maketitle

\paragraph{What these notes are.}  Suggestions to organize the ideas already in
\texttt{proposal\_eee\_draft.tex} based on (i) the five papers in
\texttt{papers/} and (ii) the empirical evidence already in
\texttt{context/}. The goal is to help you find (a) the fundamental tension
that makes the paper revealing, (b) the CF that exploits that tension, and (c)
concrete answers to the four modelling questions you flagged: unobserved
heterogeneity, learning with discrete signals, the school's role in a
Stackelberg game, and equilibrium.

\section{The fundamental tension and what it should buy you}

\paragraph{Where the proposal currently sits.}  The current draft mixes two
ideas: (1) students learn about ability from grades and recommendations, (2)
schools may distort recommendations strategically. Each is interesting, but
neither alone is a top-5 contribution: (1) is well-mapped by AAMR (2025) and
Pastorino (2024), and (2) without a sharp identification argument is
indistinguishable from descriptive school heterogeneity.

\paragraph{The tension that the data already gives you.}  The
\texttt{update\_april\_2026} memo identifies it without naming it: \emph{the
incentive of a school to use its informational advantage strategically depends
on whether the school keeps the student}. Cycle-1-only schools (41\% of the
sample) lose the student mechanically at the end of grade 2. Cycle-1+2+3
schools keep the student and face capacity constraints, pass-rate concerns,
and track-composition incentives. Within the same students-types pool, the
recommendation behaviour of these two school architectures should differ
\emph{only because of incentives}. This is the Pastorino (2024) ``differential
informativeness of jobs'' argument relabeled for schools, and it is the cleanest
identification source you have.

\paragraph{Sharpening it.}  In \texttt{update\_08\_12} the descriptive evidence
shows that, conditional on standardized exam grades, IQ, and effort, teachers
still rely heavily on \emph{relative class rank} when issuing horizontal
(STEM/non-STEM) advice. Relative rank is exactly the object that a strategic
school manipulates if it is capacity-constrained. So the tension is not just
``schools may be biased'': it is that \emph{the very statistic teachers use
(class rank) is endogenous to the school's allocation problem}, and this
distorts the signal value of the recommendation for the student's posterior
about her own track-specific ability.

\paragraph{The revealing CF.}  Once that tension is the engine of the model,
the natural CF is not ``ban tracking'' (Fu--Mehta) or ``move test cutoffs''
(Larroucau--Rios). It is:
\begin{enumerate}[leftmargin=*]
  \item \textbf{Information policy:} replace the school-issued recommendation
  with a centrally-issued, capacity-free signal of equal informativeness about
  $\theta$ (e.g., a standardized tracking score). What share of mismatch
  (downstream switching/dropout/HE field of choice) is due to schools'
  strategic distortion, vs.\ to genuine informational frictions in the
  student's own learning?

  \item \textbf{Decoupling supply from advice:} hold school capacities fixed
  but require schools to issue advice as if they had no retention incentive
  (the ``cycle-1 benchmark'').  This isolates the welfare cost of the
  Stackelberg leadership of the school.

  \item \textbf{Stackes-of-information rebalancing:} make recommendations
  binding (as B/C certificates already are) vs.\ purely advisory, and study
  who gains and who loses by SES, ability and gender. This connects to
  Larroucau--Rios's mechanism-design CFs in a way that uses your strategic
  sender as the policy lever rather than the application algorithm.
\end{enumerate}
Of these, (1) is the \emph{revealing} one because it directly decomposes the
welfare effect of recommendations into ``information'' and ``strategic
distortion'' components. The other two are robustness/design exercises.

\paragraph{Why this would be top-5.}  It is the first paper that estimates a
Bayesian-persuasion-style information design model in a real schooling
setting, with continuous and discrete signals jointly, with a credible
identification of the strategic margin (cycle-length variation), and where the
counterfactual answers a question that none of the related papers can: what is
the welfare value of \emph{the strategic component} of teacher advice? Fu--Mehta
quantify the value of tracking, Pastorino quantifies the value of learning,
AAMR quantifies the value of imperfect information; you would quantify the
value of \emph{strategic} information provision in education.

\section{Operationalizing unobserved heterogeneity ($\theta^k, \theta^u$)}

\paragraph{What you can identify from your measurement system.}  Your data is
unusually rich. Following AAMR (2025) and Cunha--Heckman--Schennach (2010), a
factor model in baseline measurements pins down the joint distribution of the
\emph{initially-known} types $\theta^k$:
\begin{itemize}[leftmargin=*]
  \item \textbf{Cognitive:} Verbal IQ, Numeric IQ, Spatial IQ at start grade 1
  $\Rightarrow$ at least two factors $(\theta^k_{\text{num}},
  \theta^k_{\text{verb}})$ identified from intra-test correlations.
  \item \textbf{Non-cognitive:} performance motivation, test anxiety,
  conscientiousness, agreeableness, extraversion (Big-5 already estimable
  from the well-being/self-concept battery) $\Rightarrow$ a non-cog factor
  $\theta^k_{\text{noncog}}$ (your update\_08\_12 already shows
  conscientiousness is independently predictive of advice).
  \item \textbf{Parental press / SES:} the parental questionnaire (education,
  occupational situation, educational climate) $\Rightarrow$ a fourth factor
  $\theta^k_{\text{parental}}$ that the descriptive evidence shows matters
  for sorting at grade 1 but is largely absorbed by school FE for advice ---
  i.e., its main role is in the \emph{first-stage} school choice, not in
  recommendation.
\end{itemize}
Three points that follow:
\begin{enumerate}[leftmargin=*]
  \item Your draft says ``$\theta^k$ has 4 elements''. The data supports
  4--5: numeric, verbal, non-cog, parental press, and arguably a separate
  motivation/anxiety dimension. Test the rank of the measurement system before
  fixing the dimension.
  \item Initial sorting $F(\theta^k\mid s)$ is non-trivial
  (\texttt{update\_april\_2026}). You have two clean handles for it: (a)
  recover $\theta^k$ from baseline measurements, condition on it; (b) use
  distance-to-school as an excluded shifter in a joint first-stage school
  choice model. Your memo proposes both. I would do (a) for the working paper
  and add (b) as a robustness exercise: distance is exogenous to the
  recommendation rule conditional on $\theta^k$, but residential location is
  endogenous in the long run, so combining the two makes the argument
  airtight.
  \item \textbf{What is $\theta^u$ here?} In AAMR/Pastorino it is fundamental
  ability. In your setting the natural object is \emph{track-specific match
  quality}: the realization the student would experience if assigned to ASO
  vs.\ TSO vs.\ BSO, and within ASO/TSO to STEM vs.\ non-STEM. The student
  has a prior over the vector of match qualities; grades update beliefs
  about each component (because cognitive skills predict track-specific
  outcomes differently); recommendations add discrete information. This is
  the AAMR multivariate correlated-learning structure applied to your tracks.
\end{enumerate}

\paragraph{Practical recipe.}  Build the measurement system at $t=1$ to
recover $F(\theta^k)$ via mixtures-of-normals (Bunting--Diegert--Maurel 2025
allows you to do this without elicited beliefs and without assuming the
measurements are selection-free, which is convenient because your tests
\emph{are} given to all students). Use a low number of discrete types $I=3$ or
$I=4$ (Pastorino) for tractability if a continuous factor model is too
expensive in the dynamic estimation. You then treat $\theta^k$ as known to
agent and school but unknown to the econometrician, and let the unknown
$\theta^u$ enter only through track-specific match quality, where learning
operates.

\section{Modelling learning with discrete signals}

\paragraph{Continuous + discrete signal structure.}  At each end-of-year
$t \in \{1,2,4,6\}$ you observe two continuous signals (Math and Dutch test
scores) and two discrete signals (the teacher's choice advice, and at grades
2/4/6 the A/B/C certificate). Following AAMR/BDM, write
\[
  G_{ijt} = \mu_j(\theta^k_i, X_{it}) + \lambda_j \theta^u_{ij} + \varepsilon_{ijt},
  \qquad j \in \{\text{Math, Dutch}\}, \quad
  \theta^u_i \sim \mathcal{N}(0,\Sigma_u),
\]
with $\Sigma_u$ allowed to be non-diagonal (correlated learning across track
dimensions, as in AAMR). Posterior on $\theta^u$ is Gaussian conjugate.

\paragraph{Adding the discrete recommendation.}  The recommendation
$R_{it}\in\mathcal{R}$ that the school issues is informative about $\theta^u$
\emph{only} insofar as the school uses additional information (its own
observations of effort, peer comparisons, knowledge of the cohort
distribution).  Two options for the student's updating, in increasing order
of complexity:
\begin{itemize}[leftmargin=*]
  \item \textbf{Reduced-form signal extraction.}  Treat the recommendation as
  a noisy categorical signal of $\theta^u$ with school-type-specific
  informativeness $\alpha_s, \beta_s$ (Pastorino's $\alpha_{jk},\beta_{jk}$).
  The student's posterior depends on the school-type-specific likelihood
  $\Pr(R_{it}=r\mid \theta^u, s)$.  Since the student observes $s$, she
  knows the likelihood she should use. \emph{This is where strategic
  scepticism enters:} a Bayesian student facing a school she suspects is
  strategic uses a flatter (less informative) likelihood. Heterogeneity in
  trust can be modelled as random coefficient in $\alpha_s$.

  \item \textbf{Structural signal extraction with strategic school.}  The
  student knows the school's recommendation rule
  $R_{it} = \rho_s(\theta^k, G_{it}, \text{cohort rank}, \text{capacity})$ and
  inverts it.  Then $R_{it}$ is informative about $\theta^u$ both directly
  (because the school sees more) and indirectly (because the school's
  own decision rule reflects what the school's information was). This is
  the Bayesian-persuasion / strategic-sender layer.
\end{itemize}
The first is enough for a first paper. The second is the version that
delivers the ``strategic information design'' contribution. I would estimate
both and use the difference as the main quantitative result.

\paragraph{Identification of the discrete-signal channel.}  Critical
observation: in your data there are cases where Math grade and Dutch grade
disagree (high Math, low Dutch or vice versa) and the recommendation
sometimes goes one way, sometimes the other. This within-student variation
in the alignment of continuous signals identifies how the school weights
each signal and, residually, what \emph{additional} information the school is
using (effort, behaviour, capacity, family pressure). The \texttt{update\_08\_12}
finding that ``even after conditioning on IQ, std grades and effort, teachers
still use relative class rank'' is directly this residual: it isolates the
school's private information component, which is the input to the strategic
distortion.

\paragraph{Pastorino-style update formula.}  The expression in your draft,
\[
  P^{s}\mid \theta = \frac{p(\theta)\, p(r\mid \theta)}{\sum_{r\notin R}p(r\mid \theta)},
\]
should be rewritten as a Bayesian update of the joint posterior on
$(\theta^u_{\text{ASO}}, \theta^u_{\text{TSO}}, \theta^u_{\text{BSO}},
\theta^u_{\text{STEM}}, \theta^u_{\text{nSTEM}})$:
\[
  f(\theta^u\mid G_t, R_t, s) \propto f(\theta^u\mid G_t)\cdot
  \Pr(R_t\mid \theta^u, s),
\]
with the first factor coming from the Gaussian update on grades and the second
the discrete-signal likelihood. The school index $s$ enters the second factor
because the recommendation rule is school-type-specific.  This makes
explicit that grades and recommendation are \emph{correlated} signals, and the
identification of $\Sigma_u$ comes from the joint distribution of grades and
recommendations across years.

\section{The school's problem and the Stackelberg game}

\paragraph{Schools as Fu--Mehta with information design.}  Fu--Mehta (2017)
gives you the right scaffolding: schools choose a tracking regime (in your
case, recommendation policy) and effort, parents (or students) respond. The
piece you add is that in your setting the school's lever is not (only) effort
or assignment but the \emph{informativeness} of a discrete signal that
students use to update beliefs.

A workable specification:
\begin{itemize}[leftmargin=*]
  \item \textbf{State of the school:} cohort distribution of types
  $G_s(\theta^k)$ inherited from sorting at grade 1; capacity constraints
  $\bar{n}_{\text{ASO-stem}}, \bar{n}_{\text{ASO-nstem}}, \dots$
  (cycle-length and supply-driven --- you observe which subtracks each school
  offers).
  \item \textbf{Action:} a recommendation rule
  $\rho_s : (\theta^k, G_t, \text{rank}, \text{effort}) \mapsto \Delta(\mathcal{R})$
  --- a (possibly mixed) mapping from observables to a probability
  distribution over recommendations $\mathcal{R} = \{\text{ASO-stem,
  ASO-nstem, TSO-stem, TSO-nstem, BSO, repeat (B), fail (C)}\}$.
  \item \textbf{Objective:} weighted sum of (a) student average final
  outcome (test scores at grade 6, HE entry, labour-market outcome),
  (b) share of students above proficiency (the A-certificate share is
  observable in your data and is the natural Fu--Mehta proficiency target),
  (c) cost of having to repeat / over-stretch capacity. Cycle-1-only schools
  put weight $\approx 0$ on (b) and (c) for grades $\geq 3$ because the
  student leaves --- this is the source of identification.
  \item \textbf{Constraint:} capacity constraints; centrally-set rules for
  who can be assigned to which track (B/C certificate rules).
\end{itemize}

\paragraph{Timing (Stackelberg).}
\begin{enumerate}[leftmargin=*,label=Stage \arabic*.]
  \item School commits to recommendation rule $\rho_s$ (chosen once per
  cohort, plausibly stable across cohorts within a school --- testable using
  the multiple cohorts in LOSO).
  \item Students observe school type $s$ and choose school at grade 1,
  knowing distance, peers, $\rho_s$ (or only the school type, if you take
  the empirical finding that A-certificate rates do not predict choice as
  evidence of myopic enrollment --- this is the simpler benchmark).
  \item Period-by-period: student exerts effort, receives grade, school
  observes everything and issues recommendation according to $\rho_s$,
  student updates beliefs.
  \item At grade 2 (the binding moment): school issues certificate +
  recommendation; student chooses track conditional on it.
  \item Continue dynamically until grade 6 / HE entry / LM entry.
\end{enumerate}

\paragraph{Equilibrium.}  An equilibrium is a tuple
$\big(\rho_s^*\big)_{s \in \mathcal{S}},\ \pi^*\big)$ where:
\begin{enumerate}[leftmargin=*]
  \item For each school type $s$, $\rho_s^*$ maximizes the school's objective
  given the induced student best-response $\pi^*$.
  \item For each student type, $\pi^*(\theta^k, s, G_t, R_t)$ is the
  Bayesian-Nash dynamic discrete choice over tracks and effort, given
  posterior beliefs over $\theta^u$ that are computed using $\rho_s^*$ as
  the likelihood of the recommendation signal.
  \item The cohort-level distribution of types and outcomes implied by
  $(\rho_s^*, \pi^*)$ is consistent with the capacity constraints and
  proficiency objective fed back into the school's problem.
\end{enumerate}
This is a perfect Bayesian equilibrium of an information-design game with
a population of receivers, in the spirit of Kamenica--Gentzkow with
heterogeneous priors and an outside option.

\paragraph{Static vs.\ dynamic.}  Your draft asks ``static or dynamic''. My
recommendation: \emph{static for the school, dynamic for the student}. A
school's recommendation rule is plausibly committed for the cohort and stable
across cohorts (you can test this with cross-cohort variation). A student is
genuinely forward-looking: choosing ASO-stem at grade 3 has option value over
HE STEM. So model the student as a finite-horizon dynamic discrete choice
agent (AAMR-style) and the school as a one-shot designer of the
recommendation rule. This is much more tractable than a fully dynamic
school problem and probably loses very little.

\section{Identification: where each ingredient comes from}

\begin{itemize}[leftmargin=*]
  \item \textbf{$F(\theta^k)$ and $F(\theta^k\mid s)$:} factor model on
  baseline measurements (BDM 2025, Cunha--Heckman--Schennach 2010).
  Initial sorting is a separate object, identified using distance.

  \item \textbf{Grade equation parameters and $\Sigma_u$:} BDM/AAMR semiparametric ID
  from the panel of grades and choices. You do not need elicited beliefs
  because BDM (2025) shows the model is identified without them under
  normality. This is the strongest part of your data.

  \item \textbf{Recommendation rule $\rho_s$:} reduced-form multinomial logit
  on observables, with school-type fixed effects, gives you the descriptive
  $\rho_s$. The structural piece needs the equilibrium condition that the
  school is optimizing given student best-response. The cycle-1 vs.\
  cycle-1+2+3 contrast plays the role of an excluded variation that shifts
  the school's incentives but not the student's primitives.

  \item \textbf{Student trust / strategic scepticism:} identified from
  the gap between how strongly students react to a recommendation in
  cycle-1 schools (where the school is presumed truthful) vs.\ in
  cycle-1+2+3 schools (where the school is suspected of distortion),
  conditional on $\theta^k$. Heterogeneity in trust by SES is identified
  from the same source.

  \item \textbf{Welfare and CFs:} once $\rho_s^*$ is identified for each
  school type, simulate the student dynamic problem under (a) the actual
  $\rho_s^*$, (b) the cycle-1 (presumed truthful) $\rho_s^*$ applied
  uniformly, (c) a centrally-set, common, capacity-free signal of the same
  Blackwell informativeness as the actual recommendations, and compute
  outcome distributions.
\end{itemize}

\section{Concrete edits I would suggest to the proposal}

\begin{enumerate}[leftmargin=*]
  \item Frame the contribution as ``estimating a strategic-information
  design model in education'' rather than ``adding a learning model with
  schools''. Cite Kamenica--Gentzkow, Bergemann--Morris, and frame
  schools as senders in the spirit of Pastorino's two-sided learning.

  \item Replace the open ``static vs.\ dynamic'' question by the explicit
  static-school / dynamic-student split.

  \item Make the cycle-1 vs.\ cycle-1+2+3 contrast the central
  identification source for the strategic margin, and announce it in the
  introduction. This is what makes the paper feel ``revealing'' rather
  than descriptive.

  \item Promote the multivariate correlated-learning structure (one
  $\theta^u$ per track-subtrack combination) explicitly. AAMR is the
  template; your tracks are 5, theirs are also 5.

  \item Drop peer effects from the model.  They are interesting but the
  capacity-constraint argument already captures the main supply-side
  force, and adding peer effects on top will create an
  endogeneity/identification problem that is not central.

  \item Drop effort as a separate model object on the student side. As you
  rightly note in the draft, it is captured in learning. Keep effort only
  as an input the school observes (it is in your data) so that the
  recommendation rule can depend on it. This matches \texttt{update\_08\_12}
  which finds effort matters for vertical (A/B/C) but not horizontal advice.

  \item In the CF section, lead with the centrally-issued signal (CF 1
  above) as the main result, and use the other CFs as robustness.
\end{enumerate}

\section{Risks and open questions}

\begin{itemize}[leftmargin=*]
  \item \textbf{Are cycle-1 schools really non-strategic?}  They have no
  retention incentive but they may still care about reputation
  (placement of their students into good cycle-2 schools). Test by looking
  at whether cycle-1 schools whose graduates feed into specific cycle-2
  schools have systematically different recommendations.
  \item \textbf{Identification of $\theta^u$ vs.\ time-varying shocks.}
  AAMR/BDM both impose normality. Verify that mixture-of-normals fits the
  empirical distribution of grades by track at $t=1$ before locking in.
  \item \textbf{Stability of $\rho_s$ across cohorts.}  The model's
  Stackelberg structure assumes school commits.  Test using LOSO's
  multi-cohort follow-ups whether $\rho_s$ moves with cohort
  composition.  If it does, you may need a school problem with state
  rather than a one-shot design.
\end{itemize}

\end{document}
